\chapter{Distributed Model-Predictive Control}
\label{ch:dmpc}

Each agent controls its motion with the model-predictive controller of
\cref{sec:mpc}. This chapter couples those controllers into a \emph{distributed}
controller (DMPC) that drives the cluster to a coordinated objective---a
formation, a trajectory, a coverage pattern---while enforcing the safety-critical
inter-agent constraints, using only neighbour-to-neighbour communication.

\section{Cluster objective and coupling}

Let each agent have a desired offset $\vv{d}_i$ from a common formation centre
$\vv{c}(t)$, which may itself follow a mission trajectory. The cluster objective
is the sum of the agents' tracking costs subject to their individual dynamics and
constraints \emph{and} to the coupling constraints that make the group
behave as one:
\begin{subequations}
\label{eq:dmpc}
\begin{align}
\min_{\{\ctrl_i\}}\ \ &\sum_{i\in\mathcal V}J_i(\state_i,\ctrl_i;\ \vv{c}(t)+\vv{d}_i)\\
\text{s.t.}\ \ &\state_{i,k+1}=\vv{f}_{d,i}(\state_{i,k},\ctrl_{i,k}),\quad
\vv{c}_i(\state_{i,k},\ctrl_{i,k})\le\vv{0}, && \forall i,k,\\
&\norm{\vv{p}_i-\vv{p}_j}\ge D, && \forall (i,j),\ i\ne j,
\label{eq:collision}
\end{align}
\end{subequations}
where $J_i$ is agent $i$'s MPC cost with reference slot $\vv{c}(t)+\vv{d}_i$, and
\cref{eq:collision} is the pairwise minimum-separation (collision-avoidance)
constraint with safety radius $D$. Two features make \cref{eq:dmpc} a genuinely
distributed problem: the reference $\vv{c}(t)$ must be agreed upon without a
central node, and the coupling constraint \cref{eq:collision} links pairs of
agents that must nonetheless decide their inputs locally.

\section{Consensus on the formation reference}

The agents agree on the formation centre by consensus. Each agent holds a local
estimate $\vv{c}_i$ of the centre (from its share of the mission command or its
own sensing) and runs, at each control step, a few rounds of the Metropolis-%
weighted iteration of \cref{eq:metropolis},
\begin{equation}
\vv{c}_i\leftarrow W_{ii}\vv{c}_i+\sum_{j\in\Neigh_i}W_{ij}\vv{c}_j,
\label{eq:centre-consensus}
\end{equation}
which converges to a common centre for any connected graph
(\cref{sec:graph}). Its rate is set by the algebraic connectivity
$\lambda_2(\Lap)$; a well-connected cluster agrees in a handful of rounds, well
within one control period. Agent $i$ then sets its MPC reference slot to
$\vv{c}_i+\vv{d}_i$ and calls \texttt{mpc\_set\_reference}
(\cref{sec:capi}). The desired offsets $\vv{d}_i$ define the formation
\emph{shape}; changing them---on command or in response to a fault---reconfigures
the formation (\cref{ch:formation}).

\section{Coupling constraints by decomposition}

The collision-avoidance coupling \cref{eq:collision} links neighbouring agents'
decision variables. There are two established routes to solving the coupled
problem \cref{eq:dmpc} with only local computation and neighbour messages, and
the framework supports both.

\paragraph{Dual decomposition / ADMM.} Introduce a copy of each shared quantity
and a consistency constraint tying an agent's copy to its neighbours'; the
augmented Lagrangian of \cref{eq:dmpc} then \emph{separates} across agents given
the dual variables, and the alternating-direction method of multipliers (ADMM)
alternates a local solve---one call to each agent's MPC---with a neighbour
exchange and a dual update \cite{boyd2011admm,summers2012distributed}. Each ADMM
iteration is one local MPC solve plus one round of communication, so a fixed
(small) number of iterations per control step keeps the scheme real-time; the
per-agent solve is exactly the RTI step of \cref{sec:mpc}, warm-started from the
previous iterate. Accelerated dual-decomposition variants
\cite{giselsson2013accelerated} reduce the iteration count. This is the
\emph{distributed-optimisation} route: at convergence it solves the centralised
problem, and terminated early it yields a feasible, if suboptimal, coordinated
input.

\paragraph{Cooperative DMPC and its guarantees.} When the agents optimise the
\emph{plant-wide} objective rather than only their own, the cooperative-DMPC
theory of Stewart \etal \cite{stewart2010cooperative} applies: the systemwide
cost is a common Lyapunov function, so the closed loop is stable and recursively
feasible \emph{at every intermediate iterate}, including a single iteration per
sample, provided each agent warm-starts from the previous feasible plan. This
``any-iterate'' property is exactly what a real-time cluster needs, and the
distributed synthesis of terminal ingredients of Conte \etal
\cite{conte2016distributed} certifies stability with only neighbour-to-neighbour
coupling. We state these as the guarantees the architecture inherits; the review
of \cite{christofides2013distributed} and the text of \cite{rawlings2017model}
develop the cooperative and non-cooperative variants and their proofs.

\section{Real-time safety filter}
\label{sec:cbf}

For hard, instantaneous collision avoidance the framework layers a
control-barrier-function (CBF) \emph{safety filter} between the MPC command and
the actuator, which minimally alters the commanded acceleration to keep the pair
separation constraint invariant \cite{ames2017control,wang2017safety}. Define the
pairwise barrier
\begin{equation}
h_{ij}(\vv{p})=\norm{\vv{p}_i-\vv{p}_j}^2-D^2\ \ge0,
\label{eq:cbf-h}
\end{equation}
whose safe set is exactly \cref{eq:collision}. Because \cref{eq:cbf-h} has
relative degree two with respect to the commanded acceleration (a double
integrator in position), we use the higher-order CBF condition
\begin{equation}
\ddot h_{ij}+k_1\dot h_{ij}+k_2 h_{ij}\ge0,
\label{eq:hocbf}
\end{equation}
with $k_1,k_2>0$ chosen so the associated characteristic polynomial is Hurwitz.
Writing $\Delta\vv{p}=\vv{p}_i-\vv{p}_j$, $\Delta\vv{v}=\vv{v}_i-\vv{v}_j$ and
conservatively treating the neighbour's acceleration as zero over the step (the
standard decentralised assumption, since each agent actuates only its own input),
$\ddot h_{ij}=2\,\Delta\vv{v}\T\Delta\vv{v}+2\,\Delta\vv{p}\T\vv{a}_i$ and
\cref{eq:hocbf} becomes the \emph{linear} inequality in agent $i$'s acceleration
$\vv{a}_i$
\begin{equation}
-2\,\Delta\vv{p}\T\vv{a}_i\ \le\ 2\,\Delta\vv{v}\T\Delta\vv{v}+k_1\dot h_{ij}+k_2 h_{ij}.
\label{eq:cbf-row}
\end{equation}
(If the neighbour's acceleration $\vv{a}_j$ is communicated, the exact row carries
the additional term $-2\,\Delta\vv{p}\T\vv{a}_j$ on the right-hand side; treating
$\vv{a}_j=\vv{0}$ is conservative and needs no communication.)
The safe acceleration is the solution of the small quadratic program that stays
closest to the MPC command $\vv{a}_i^{\text{cmd}}$ subject to one row
\cref{eq:cbf-row} per neighbour,
\begin{equation}
\vv{a}_i^{\text{safe}}=\arg\min_{\vv{a}_i}\ \norm{\vv{a}_i-\vv{a}_i^{\text{cmd}}}^2
\quad\text{s.t. \cref{eq:cbf-row}}\ \forall j\in\Neigh_i,
\label{eq:cbf-qp}
\end{equation}
a problem of the agent's control dimension that is solved in microseconds by the
same constraint machinery as the core solver. The filter is inactive when the
MPC command is already safe and engages only near the safety boundary, giving a
provably forward-invariant safe set under the standard CBF assumptions
\cite{ames2017control}. The collision-avoidance rows may equivalently be written
directly into the MPC constraint $\vv{c}_i$ of \cref{eq:ocp}; the safety filter is
the lightweight, always-feasible last line of defence that the simulation study
uses.

\begin{algorithm}[t]
\caption{Distributed MPC step at agent $i$}
\label{alg:dmpc}
\DontPrintSemicolon
\KwIn{state estimate $\est{\state}_i$ (from \cref{alg:dmhe}); local centre
estimate $\vv{c}_i$; neighbour states/centres}
consensus rounds on the centre: \cref{eq:centre-consensus}\;
set reference slot $\vv{c}_i+\vv{d}_i$; \texttt{mpc\_prepare} (matrix pass)\;
$\vv{a}_i^{\text{cmd}}\leftarrow$ \texttt{mpc\_feedback}$(\est{\state}_i)$\hfill// RTI, \cref{eq:rti-feedback}\;
$\vv{a}_i^{\text{safe}}\leftarrow$ CBF-QP \cref{eq:cbf-qp} over $\Neigh_i$\;
differential flatness \cref{eq:flat1}--\eqref{eq:flat2} + inner attitude loop
$\rightarrow$ actuators\;
broadcast $(\vv{p}_i,\vv{v}_i)$ to $\Neigh_i$
\end{algorithm}
